\section{From stable end profiles to the monotone theorem}
\label{sec:monotone}

In this part we deduce Theorem~\ref{thm:monotone} from the unrestricted
seven-dimensional classification in Theorem~\ref{thm:stable}, whose
independent analytical proof occupies Parts II--IV.  None of those parts
uses the geometric deduction given here.  The end
profiles will initially be allowed to be planar heteroclinics.  Their
minimizing property suffices for the argument; no well limit is imposed on
the given solution.  We first derive minimization and surface-order energy
growth, then identify the complete limiting energy with perimeter.  Only
after this exact identification do we use the geometry of a directed
minimizing cone and return to the equation at unit density.

We use the energy $\mathcal E$ and global minimizers defined in
Section~\ref{sec:introduction}.  It is enough to consider
competitors valued in $[-1,1]$: truncation to this interval decreases both
terms of the energy.  This definition also gives minimization in every
bounded Lipschitz domain: extend a same-trace competitor by $u$ to a
containing ball, use the Sobolev gluing property across the boundary, and
cancel the identical exterior energies.  We retain the normalization
\begin{equation}
 s_0=\int_{-1}^1\sqrt{2W(a)}\dd a
     =\int_\R g'(s)^2\dd s
     =\frac{2\sqrt2}{3}=\frac\sigma2.
 \label{eq:monotone-tension}
\end{equation}

\subsection{Stable end profiles and a variational comparison}

The first task is to turn strict monotonicity into minimizing barriers.
The positive directional derivative gives stability, and a long-cylinder
test transfers that stability to the seven-dimensional limits.  Their
classification then supplies a calibration, even when a limit is a planar
profile rather than a constant well.

\begin{lemma}[End profiles]
\label{lem:monotone-ends}
Let $u\in C^2(\R^8,(-1,1))$ solve $\Delta u=u^3-u$ and satisfy
$u_8>0$.  There are smooth stable solutions
$v_-,v_+:\R^7\to[-1,1]$ such that
\begin{equation}
 u(x',t+s)\longrightarrow v_\pm(x')
 \quad\text{in }C^k_{\mathrm{loc}}(\R^8)
 \quad\text{as }s\to\pm\infty,
 \qquad k\ge0.
 \label{eq:monotone-end-convergence}
\end{equation}
They satisfy $v_-(x')<u(x',t)<v_+(x')$ everywhere.  Their extensions
independent of the eighth coordinate are global minimizers in $\R^8$.
\end{lemma}

\begin{proof}
Elliptic interior estimates, applied on unit balls, give smoothness and
uniform bounds for every fixed derivative order of $u$.  The function
$h=u_8>0$ satisfies $(-\Delta+3u^2-1)h=0$.  For every compactly supported
$C^1$ function $\eta$, integration by parts gives the ground-state identity
\begin{equation}
 \int_{\R^8}\bigl(|\nabla\eta|^2+(3u^2-1)\eta^2\bigr)
   =\int_{\R^8}h^2\left|\nabla\left(\frac\eta h\right)\right|^2
   \ge0.
 \label{eq:monotone-ground-state}
\end{equation}
The division is legitimate because $h$ has a positive minimum on the
support of $\eta$.

Bounded monotonicity gives pointwise limits along each vertical line.
Translation compactness and the derivative bounds improve this to
\eqref{eq:monotone-end-convergence}.  Indeed every convergent subsequence
has the same pointwise limit, and a bounded vertical shift leaves that
limit unchanged.  Passing the equation and stability to these limits first
gives solutions and stability in eight dimensions.  To see stability in
seven dimensions, fix $\varphi\in C_c^1(\R^7)$ and a nonzero
$\zeta\in C_c^1(\R)$.  Test with
$\eta(x',t)=\varphi(x')\zeta(t/L)$ and divide by
$L\int\zeta^2$.  The result is
\[
 \int_{\R^7}\bigl(|\nabla\varphi|^2+(3v_\pm^2-1)\varphi^2\bigr)
 \ge -\frac1{L^2}\frac{\int|\zeta'|^2}{\int\zeta^2}
                         \int_{\R^7}\varphi^2.
\]
Letting $L\to\infty$ proves the assertion.  For example,
$u(x',t)<u(x',t+1)\le v_+(x')$, which proves the strict upper
inequality; the lower one is analogous.

Theorem~\ref{thm:stable} says that each $v_\pm$ is a constant well or
a planar heteroclinic.  To check the minimizing assertion directly, put
$\Phi(a)=\int_{-1}^a\sqrt{2W(b)}\dd b$.  For a unit vector $e$ and a
competitor $w\in H^1(\Omega,[-1,1])$,
\begin{equation}
 \frac12|\nabla w|^2+W(w)
    \ge \sqrt{2W(w)}\,\partial_e w
    =\partial_e\Phi(w).
 \label{eq:monotone-calibration}
\end{equation}
There is equality when $w(x)=g(e\cdot x-c)$, since
$g'=\sqrt{2W(g)}$.  The integral of the last expression depends only
on the boundary trace.  Thus such a profile minimizes on every ball, in
any ambient dimension in which $e$ is a unit vector.  This includes the
extensions of the seven-dimensional profiles, with $e_8=0$.  The
constants $\pm1$ have zero energy and also minimize.
\end{proof}

The next comparison also applies in other dimensions.  It records exactly
which property of the end profiles is used.  The constrained comparison
between end profiles was proved by Alberti, Ambrosio and Cabr\'e
\cite[Theorem~4.4]{AlbertiAmbrosioCabre2001}.  The unrestricted
minimization implication is the classical theorem of
\citet[Theorem~1.3, p.~441; proof on pp.~452--453]{JerisonMonneau2004};
we include the compact comparison argument in the present normalization.
Its globally $C^{1,1}$ potential hypothesis is met by
\[
 \widetilde W(s)=
 \begin{cases}
  (s+1)^2,&s<-1,\\
  W(s),&-1\le s\le1,\\
  (s-1)^2,&s>1.
 \end{cases}
\]
The values and the first two derivatives match at the wells, and
$\widetilde W'$ is globally Lipschitz.  Both potentials are nonnegative
and agree on $[-1,1]$.  Clipping any competitor to this interval decreases
its gradient energy and replaces its exterior potential by zero for
either potential, so their minimizing comparisons for $u$ agree.

\begin{lemma}[Minimizing barriers and compact sliding]
\label{lem:monotone-minimizer}
Let $N\ge2$, and suppose a bounded solution
$u:\R^{N-1}\times\R\to(-1,1)$ of $\Delta u=W'(u)$ is strictly
increasing in the last coordinate.  Suppose its vertical limits
$a_\pm(x')$, extended to $\R^N$, are global minimizers.  Then $u$ is a
global minimizer.  Moreover there is a dimensional constant $C_N$ such
that
\begin{equation}
 \mathcal E(u;B_R(y))\le C_N R^{N-1}
 \qquad (y\in\R^N,\ R\ge2).
 \label{eq:monotone-energy-growth}
\end{equation}
\end{lemma}

\begin{proof}
Fix a ball $B$.  The direct method produces an energy minimizer $w$ with
boundary trace $u$; truncation allows $|w|\le1$.  First clip it from
below by $a_-$, writing $w_1=\max\{w,a_-\}$.  The Sobolev lattice
identity is
\[
 \mathcal E(\max\{w,a_-\};B)
 +\mathcal E(\min\{w,a_-\};B)
 =\mathcal E(w;B)+\mathcal E(a_-;B).
\]
Since $u>a_-$ on $\partial B$, the minimum on the left has the same
trace as $a_-$.  Minimality of $a_-$ therefore shows
$\mathcal E(w_1;B)\le\mathcal E(w;B)$.  The maximum has trace $u$,
so it remains an unconstrained minimizer.  The analogous upper clipping
gives another unconstrained minimizer
$\widetilde w=\min\{w_1,a_+\}$ satisfying
$a_-\le\widetilde w\le a_+$.

The Euler--Lagrange equation and elliptic boundary regularity make
$\widetilde w$ smooth up to $\partial B$.  A difference of two
solutions $b,c$ here satisfies
$-\Delta(b-c)+(b^2+bc+c^2-1)(b-c)=0$, with bounded zeroth-order
coefficient.  Strong comparison for a nonnegative solution of this
linear equation, together with the strict boundary inequalities, gives
\[
 a_-<\widetilde w<a_+\quad\text{on }\overline B.
\]
These gaps have positive minima on this compact set.  Vertical translates
$u_s(x',t)=u(x',t+s)$ converge uniformly there to $a_+$ as
$s\to+\infty$, so $u_s>\widetilde w$ for all sufficiently large $s$.
Let
\[
 s_* =\inf\{s\ge0:u_s\ge\widetilde w\text{ on }\overline B\}.
\]
The set inside the infimum is an upper interval, by strict monotonicity.
Continuity gives $u_{s_*}\ge\widetilde w$.  If $s_*>0$, then
$u_{s_*}>u=\widetilde w$ on $\partial B$, and strong comparison gives
strict inequality throughout $\overline B$.  Compactness and continuity
then allow a smaller positive translation, a contradiction.  Thus
$\widetilde w\le u$.  For the reverse inequality, $u_{-s}$ lies below
$\widetilde w$ when $s$ is large.  Decrease $s\ge0$ until first
contact.  A positive first-contact parameter is again impossible, because
$u_{-s}<u=\widetilde w$ on $\partial B$ and strong comparison gives
a strict compact gap.  At $s=0$ this yields $u\le\widetilde w$.
Consequently $u=\widetilde w$ on $B$, and $B$ was arbitrary.

Finally choose a smooth cutoff $\eta_R$ that is zero on $B_{R-1}(y)$,
one near $\partial B_R(y)$, and satisfies $|\nabla\eta_R|\le C$.
The competitor $1+\eta_R(u-1)$ has zero energy in the inner ball and
bounded energy density in the remaining shell.  The latter assertion
uses $|u|\le1$ and the uniform unit-scale interior gradient estimate for
the equation.  The shell has volume at most $C_NR^{N-1}$, proving
\eqref{eq:monotone-energy-growth}.
\end{proof}

Lemmas~\ref{lem:monotone-ends} and~\ref{lem:monotone-minimizer} apply to
the solution of Theorem~\ref{thm:monotone}.  They establish minimization
and the required energy growth directly from its hypotheses and
Theorem~\ref{thm:stable}.

\subsection{Exact convergence to a minimizing boundary}

We have obtained the variational comparison needed for the blow-down.
The remaining issue in this step is equality of measures: phase
convergence and a perimeter lower bound alone would permit additional
diffuse energy.  We eliminate that possibility by comparing with a
recovery sequence patched to the original solution in a thin annular
layer.  This argument does not require a prescribed trace for the
recovery sequence.

We use the following established variational and geometric inputs.
For a finite-perimeter set $F$ in a bounded Lipschitz domain $\Omega$,
the Modica--Mortola recovery theorem provides, for any sequence
$\eps_j\downarrow0$, functions
$w_j\in H^1(\Omega,[-1,1])$ such that
\begin{equation}
 \begin{split}
 w_j&\longrightarrow 2\mathbf1_F-1\quad\text{in }L^1(\Omega),\\
 \limsup_j\int_\Omega
   \left(\frac{\eps_j}2|\nabla w_j|^2+\frac{W(w_j)}{\eps_j}\right)
   &\le s_0 P(F;\Omega).
 \end{split}
 \label{eq:monotone-recovery}
\end{equation}
Here $P(F;\Omega)=|D\mathbf1_F|(\Omega)$ is relative perimeter
\citep{ModicaMortola1977}; the finite-perimeter recovery statement is
also proved in \citet[Section~3, pp.~140--141]{Modica1987}.
For a locally perimeter-minimizing set in
$\R^N$, its perimeter measure defines a stationary integral varifold
with multiplicity one on the reduced boundary.  Its support is a smooth
embedded minimal hypersurface outside a closed singular set of Hausdorff
dimension at most $N-8$ (and the singular set is empty when $N\le7$).
The centered area ratio is nondecreasing, and equality at all scales
forces a cone.  These are the classical regularity, dimension-reduction,
and monotonicity results for minimizing boundaries
\citep{Federer1969,Giusti1984,Simons1968}; for the sharp singular-set
dimension bound, see \citet[Theorem~1, p.~769]{Federer1970}.
The argument below uses
recovery and comparison to prove exact energy convergence before using
this multiplicity-one geometric theory.

\begin{proposition}[Exact perimeter limit]
\label{prop:monotone-exact-limit}
Let $u:\R^N\to[-1,1]$ be an entire global minimizer with
$\mathcal E(u;B_R)\le CR^{N-1}$ for $R\ge2$.  Given
$R_j\to\infty$, set $\eps_j=R_j^{-1}$ and $u_j(x)=u(R_jx)$.
After taking a subsequence there is a locally finite-perimeter set $E$
such that
\begin{equation}
 \begin{aligned}
 u_j&\longrightarrow 2\mathbf1_E-1
                    &&\text{in }L^1_{\mathrm{loc}}(\R^N),\\
 \left(\frac{\eps_j}2|\nabla u_j|^2+
                \frac{W(u_j)}{\eps_j}\right)\dd x
   &\rightharpoonup s_0|D\mathbf1_E|
                    &&\text{as local Radon measures}.
 \end{aligned}
 \label{eq:monotone-exact-limit}
\end{equation}
The set $E$ is locally perimeter minimizing against compact perturbations.
\end{proposition}

\begin{proof}
\emph{Compactness and the lower bound.}
Recall the phase primitive
$\Phi(a)=\int_{-1}^a\sqrt{2W(b)}\,\dd b$.
Denote the rescaled energy by $\mathcal E_j$ and its energy measure by
$\mu_j$.  Scaling the minimizing comparison shows that $u_j$ minimizes
$\mathcal E_j$ against every compact modification.  A change of variables
gives
\begin{equation}
 \mu_j(B_r)=R_j^{1-N}\mathcal E(u;B_{R_jr}).
 \label{eq:monotone-rescaled-energy}
\end{equation}
Thus these measures are locally bounded.  Moreover
\[
 |\nabla\Phi(u_j)|
   \le \frac{\eps_j}2|\nabla u_j|^2+\frac{W(u_j)}{\eps_j},
 \qquad \int_K W(u_j)\le\eps_j\mu_j(K)\longrightarrow0
\]
on each compact set $K$.  Local BV compactness, followed by a diagonal
subsequence, implies that $\Phi(u_j)$ has an $L^1_{\mathrm{loc}}$
limit.  The vanishing potential forces its values to lie in
$\{\Phi(-1),\Phi(1)\}=\{0,s_0\}$.  The inverse of $\Phi$ is
uniformly continuous on $[0,s_0]$, so this also proves the asserted
local $L^1$ convergence of $u_j$ to $2\mathbf1_E-1$.  Weak compactness
of $\mu_j$ and lower semicontinuity
of total variation give a local measure limit $\mu$ satisfying
\begin{equation}
 \mu\ge s_0|D\mathbf1_E|.
 \label{eq:monotone-perimeter-lower}
\end{equation}

\emph{A recovery comparison with the correct boundary trace.}
We prove the opposite inequality, including the boundary gluing.  Fix a
ball $B_r(y)$ and $\delta>0$.  Let $F$ be a finite-perimeter set in
$B_{r+\delta}(y)$ that agrees with $E$ on the annulus
$A=B_{r+\delta}(y)\setminus\overline{B_r(y)}$.  The case $F=E$
is permitted.  Choose a recovery sequence $w_j$ on this larger ball as
in \eqref{eq:monotone-recovery}.  Both $w_j$ and $u_j$ converge on $A$
to $2\mathbf1_E-1$.  Their boundedness implies
\begin{equation}
 \int_A|u_j-w_j|^2\le2\int_A|u_j-w_j|\longrightarrow0.
 \label{eq:monotone-annulus-mismatch}
\end{equation}

For $\eps_j<\delta/4$, partition all but a remainder of $A$ into
$L_j=\lfloor\delta/\eps_j\rfloor-1$ disjoint radial layers of
thickness $\eps_j$.  Thus $L_j\ge\delta/(2\eps_j)$, and the
outer remainder leaves a neighborhood of $\partial B_{r+\delta}(y)$
untouched.  In one such layer $S$, let $\chi$ increase from
zero to one, with $|\nabla\chi|\le C/\eps_j$, and interpolate by
$z_j=(1-\chi)w_j+\chi u_j$.  The cutoff is zero on the inner side
and one on the outer side.  On $[-1,1]$, boundedness of $W''$ and
$|W'|^2\le CW$ give, uniformly for $0\le\chi\le1$,
\[
 W((1-\chi)a+\chi b)
       \le C\bigl(W(a)+W(b)+|a-b|^2\bigr).
\]
Expanding the gradient therefore yields
\begin{equation}
 \mathcal E_j(z_j;S)
 \le C\bigl(\mathcal E_j(u_j;S)+\mathcal E_j(w_j;S)\bigr)
             +\frac C{\eps_j}\int_S|u_j-w_j|^2.
 \label{eq:monotone-strip-cost}
\end{equation}
Averaging the right-hand side over the disjoint layers selects $S_j$
with cost at most
\begin{equation}
 C\frac{\eps_j}{\delta}
     \bigl(\mathcal E_j(u_j;A)+\mathcal E_j(w_j;A)\bigr)
       +\frac C\delta\int_A|u_j-w_j|^2=o_j(1).
 \label{eq:monotone-strip-average}
\end{equation}
Here $\delta$ is fixed before $j\to\infty$; the two energies are
bounded by compactness and recovery.  In particular no rate for the
$L^1$ convergence is needed: the averaging cancels the factor
$\eps_j^{-1}$ multiplying the mismatch in a single layer.

Use $w_j$ inside $S_j$, the interpolation on $S_j$, and $u_j$
outside $S_j$.  This is an admissible compact modification of $u_j$.
Its minimizing property, and nonnegativity of the original energy in
the part of the annulus inside $S_j$, give
\[
 \mathcal E_j(u_j;B_r(y))
   \le \mathcal E_j(w_j;B_{r+\delta}(y))+o_j(1).
\]
Consequently, whenever $\mu(\partial B_r(y))=0$,
\begin{equation}
 \mu(B_r(y))\le s_0 P(F;B_{r+\delta}(y)).
 \label{eq:monotone-competitor-bound}
\end{equation}
\emph{Equality of energy and perimeter, then minimality.}
First set $F=E$.  At radii with
$\mu(\partial B_r(y))=|D\mathbf1_E|(\partial B_r(y))=0$, send
$\delta\downarrow0$ in \eqref{eq:monotone-competitor-bound}.  These
are all but countably many radii at each center.  Together with
\eqref{eq:monotone-perimeter-lower}, the resulting inequalities on
balls prove $\mu=s_0|D\mathbf1_E|$ as measures.  More explicitly,
$\mu-s_0|D\mathbf1_E|$ is already a nonnegative Radon measure by
\eqref{eq:monotone-perimeter-lower}, and it vanishes on these balls,
which form a neighborhood basis.  In particular no additional
transition-layer mass remains.

Next suppose $F\mathbin\triangle E$ is compactly contained in an
open ball.  Choose $r$ outside that compact set, inside the ball, and
with zero perimeter measure on $\partial B_r(y)$.  Then $F=E$ near
this sphere, so it is a continuity sphere for both perimeter measures.
Apply \eqref{eq:monotone-competitor-bound} and let $\delta\downarrow0$
to obtain $P(E;B_r(y))\le P(F;B_r(y))$.  The perimeter measures
coincide outside the compact modification; hence this proves local
perimeter minimality, as asserted.
\end{proof}

\subsection{The directed minimizing cone in dimension eight}

The exact limit is now available together with perimeter minimality.
We next combine energy monotonicity with vertical monotonicity: the first
makes the limiting boundary a cone, while the second fixes the sign of
its vertical normal component.  In this geometric step, dimension eight
enters when we analyze the singular set of this cone.

For our eight-dimensional solution, the gradient inequality
\eqref{eq:diffuse-modica} and monotonicity formula
\eqref{eq:finite-monotonicity} give a nondecreasing normalized energy
\[
 M_r(u)=\frac{\mathcal E(u;B_r)}{s_0\omega_7r^7}.
\]
By \eqref{eq:monotone-energy-growth} it has a finite limit
$M_\infty$.  This limit is positive, since $u$ is nonconstant and its
energy is positive on some ball.  Apply
Proposition~\ref{prop:monotone-exact-limit} to any sequence
$R_j\to\infty$.  At continuity radii, its exact convergence and
\eqref{eq:monotone-rescaled-energy} imply
\begin{equation}
 P(E;B_r)=\omega_7r^7M_\infty.
 \label{eq:monotone-cone-ratios}
\end{equation}
Squeezing between continuity radii proves this identity for every $r>0$.
Thus the stationary varifold of the minimizing boundary has constant
centered area ratio.  The equality case of monotonicity makes its support
\[
 C=\supp|D\mathbf1_E|
\]
a cone with vertex at the origin.  It is nonempty and contains the origin.
Its multiplicity is one, since its weight is the perimeter measure.

Vertical monotonicity also survives this convergence.  For each
nonnegative $\psi\in C_c^1(\R^8)$,
\[
 -\int u_j\,\partial_8\psi=\int (\partial_8u_j)\psi\ge0.
\]
The $L^1$ limit gives $\partial_8\mathbf1_E\ge0$ as a distribution
and hence as a measure.  Let $\nu_E$ denote the unit normal pointing
\emph{into} $E$ on the regular boundary, so that
$D\mathbf1_E=\nu_E|D\mathbf1_E|$.  Then
\begin{equation}
 \nu_E\cdot e_8\ge0\quad\text{on the regular part of }C.
 \label{eq:monotone-directed-cone}
\end{equation}

\begin{lemma}[Flatness of the directed cone]
\label{lem:monotone-flat-cone}
Let $E\subset\R^8$ be locally perimeter minimizing.  Suppose its
nonempty boundary support $C$ is a cone with vertex $0$ and
$\partial_8\mathbf1_E\ge0$.  Then $C$ is a hyperplane, with
perimeter multiplicity one.
\end{lemma}

\begin{proof}
This is the classical directed-cone reduction; see also
\citet[Theorem~4.5 and Corollary~4.6]{Schnuerer2003}.  We give its
short proof to make the treatment of the vertex explicit.
\emph{The cone has a smooth compact link.}
The singular-set dimension bound is at most zero.  A nonzero singular
point would, by dilation invariance, generate a whole singular ray:
each dilation is an ambient diffeomorphism preserving $C$, and therefore
preserves the distinction between smooth and singular boundary points.
This is impossible, so $\operatorname{Sing}C\subset\{0\}$.
Consequently
\[
 \Lambda=C\cap\mathbb S^7
\]
is a smooth compact embedded minimal hypersurface of $\mathbb S^7$,
possibly with several connected components.  The oriented normal is
continuous on the regular boundary.  Since tangent normal lines are
dilation invariant and the choice of side is continuous, that orientation
is constant along every regular ray: it takes values in the same
two-point set along that connected ray.  This uses conicality of the
support and does not presume conicality of $E$.  The function
$f=\nu_E\cdot e_8\ge0$ is therefore homogeneous of degree zero.

The normal component of a Euclidean translation on a minimal hypersurface
satisfies the Jacobi equation
\citep[Corollary~3.3.1, p.~74]{Simons1968}
\[
 \Delta_C f+|A_C|^2f=0.
\]
The cone metric is $\dd r^2+r^2g_\Lambda$ and
$|A_C|^2=r^{-2}|A_\Lambda|^2$, where $A_\Lambda$ is the second
fundamental form of $\Lambda$ in the sphere.  Thus
\begin{equation}
 \Delta_\Lambda f+|A_\Lambda|^2f=0.
 \label{eq:monotone-link-jacobi}
\end{equation}
We separate the case where this nonnegative function is positive on a
link component from the case where it vanishes on every component.
\emph{A positive directional normal component.}
On a connected component where $f$ is not identically zero, the strong
maximum principle gives $f>0$.  Integration of
\eqref{eq:monotone-link-jacobi} on that compact component yields
$\int|A_\Lambda|^2f=0$, so it is totally geodesic.  A connected
compact embedded totally geodesic hypersurface of the sphere is a whole
equatorial $\mathbb S^6$.  Any other component, being disjoint from
this equator, would lie in one of its open hemispheres.  This cannot
happen: the corresponding signed coordinate function $\ell$ on a
closed minimal six-dimensional hypersurface of $\mathbb S^7$ satisfies
$\Delta_\Lambda\ell+6\ell=0$, hence has integral zero, whereas in
an open hemisphere it has strict constant sign.  There are consequently
no other components, and $C$ is the hyperplane over that equator.

\emph{A vanishing directional normal component.}
It remains to handle the case $f\equiv0$ on all of $\Lambda$.
Perimeter measure gives zero mass to the possible singular point, so now
\begin{equation}
 \partial_8\mathbf1_E=0\quad\text{on all of }\R^8.
 \label{eq:monotone-vertical-invariance}
\end{equation}
A locally integrable function whose distributional derivative in one
direction vanishes is invariant under translations in that direction.
For example this follows by mollifying, applying the assertion to the
smooth functions, and passing to the local $L^1$ limit.  Thus each
vertical translation preserves $\mathbf1_E$ as an $L^1_{\rm loc}$
class, and also preserves its perimeter measure, its support $C$, and
the singular set of that support.  If $0$ were singular, every point of
the vertical axis would be singular, contradicting
$\operatorname{Sing}C\subset\{0\}$.  Therefore the origin is
regular.  A cone smooth at its vertex equals its tangent hyperplane
there: dilating toward the origin keeps the cone fixed while its smooth
local graphs converge to that tangent plane.  This proves flatness in
the remaining case.  In either case the weight is
$\cH^7\mathbin\llcorner C$, by the multiplicity-one structure of a
minimizing boundary.
\end{proof}

Applying Lemma~\ref{lem:monotone-flat-cone} to
\eqref{eq:monotone-cone-ratios} shows
\begin{equation}
 M_\infty(u)=1.
 \label{eq:monotone-unit-density}
\end{equation}
This is the point where the exact measure convergence in
Proposition~\ref{prop:monotone-exact-limit} is needed: flat support
alone would not determine the diffuse energy density.

\subsection{Return to the equation and the exact profile}

The near-unit-density theorem of Wang
\citep[Theorem~11.1, p.~3036]{Wang2017} says that, in every ambient
dimension $N$, a nonconstant bounded entire Allen--Cahn solution with
the Modica bound and
\[
 \lim_{R\to\infty}\frac{\mathcal E(u;B_R)}
                           {s_0\omega_{N-1}R^{N-1}}
            \le1+\delta_N
\]
for a sufficiently small dimensional $\delta_N>0$ is one-dimensional.
The constant wells can be included as the separate trivial alternatives;
see also \citet[Theorem~2.16]{FSS2025}.  Our solution has the required
range and Modica bound, is nonconstant by $u_8>0$, and satisfies this
energy condition by \eqref{eq:monotone-unit-density}.  Hence
$u(x)=h(e\cdot x)$ for some unit vector $e$.

For completeness we recover the profile and its orientation directly.
Since $u_8=e_8h'>0$, we have $e_8\ne0$; reversing $e$ and the
scalar coordinate if necessary gives $e_8>0$ and $h'>0$ everywhere.
The finite limits $a_\pm=\lim_{s\to\pm\infty}h(s)$ are distinct
and lie in $[-1,1]$.  Translated compactness for the ODE
$h''=h^3-h$ shows $a_\pm^3-a_\pm=0$ and $h'(s)\to0$ at both
ends.  Indeed any translated subsequential limit at an end is the
constant $a_\pm$, and the uniform derivative bounds give convergence
of its derivatives.  The first integral is
\[
 \frac12h'(s)^2-W(h(s))=\text{constant}.
\]
Taking both limits gives $W(a_-)=W(a_+)$.  Among the three stationary
roots $-1,0,1$, the only distinct pair with equal potential is
$-1,1$.  Thus $a_-=-1$, $a_+=1$, and the first integral and $h'>0$
give
\[
 h'=\frac{1-h^2}{\sqrt2},\qquad
 h(s)=\tanh\left(\frac{s-c}{\sqrt2}\right)
\]
for some $c\in\R$.  This proves Theorem~\ref{thm:monotone}, including
the condition $e_8>0$.  Conversely this formula has strict range
$(-1,1)$, satisfies $g''=g^3-g$, and obeys
$u_8=e_8g'(e\cdot x-c)>0$.

